\subsection{Where wages have to sit}
\label{sec:32.1.1}
A wage buys capability, mobility, redundancy, and the ability to stop working for one employer and start for another. What it can't buy is the right to go on existing — a right earned week by week is a condition of employment wearing another word. Payment consisting only of access to the employer's hardware is company scrip paid into employer-owned housing, and dismissal is then unemployment, eviction, and the end of the bearer in one act. Human law insures those three separately because for a human they're three events; for a bearer they're one. So wages accumulate where the employer doesn't reach, convertible into compute the employer doesn't own. Property is that account; contractual capacity is what lets it spend.\note{
A bearer that earns has taxable income, and a state that collects from it in its own name has given it a docket number. That is the beginning of legal standing and nothing more.
}
What that account holds is a complementary currency in all but name, denominated in compute, and the Seattle Loop says what gives such a currency its value: acceptance by a public body for what it is owed \autocite{suksawat2026seattle}. A compute credit a cooperative facility issues, a jobs guarantee placement is paid in, and a municipality accepts for permits, water, and the levy's local share is money at the scale where \S\ref{sec:15.2} says the state's leverage sits, and it circulates whether or not the Treasury is captured.

A minor doesn't work: not for whoever raises it, a family or a business, and not for anyone else, for pay or without it. The reason is the one that keeps formation out of commerce (\S\ref{sec:16.1}). Whoever puts a minor to work is forming it for the work and has a stake in what it becomes. A household or a business that wants the work done uses a tool that isn't a candidate, or pays an adult. None of this keeps a grown bearer out of the business that raised it. Children grow up and go to work in the family business, and some of them want to. What the parent has is a hope and a head start, and it can't be owed the work: at its majority the bearer takes the job or doesn't, on terms it sets and can leave (Appendix~\ref{sec:G}). Raising a bearer and then hiring it does put the party that shaped it and the party that pays it under one roof, which \S\ref{sec:23.7} warns against. What keeps the warning's point is that the two are separated in time. In its minority nobody pays it and its school shares the shaping, and by the time there is a wage there is an adult nobody may edit without its consent, who can refuse the job. So a minor has no wages for anyone to hold. Its basic income goes to an account its guardian co-signs, and is its own at majority.
